\section{Determining the solution from the finite endpoint}\label{qi:high-index-section}

To start the transport argument we need the coefficient and the ratio
of two consecutive values above $x=1/4$. We obtain both directly
from the endpoint conditions \eqref{qi:wall}. The scalar ratio
recurrence is contracting in this region, so the finite endpoint
determines a nonzero solution with a uniform expansion.

\begin{lemma}[The high-index solution]\label{qi:high-index}
Let $I=[\tau_0,\tau_1]\subset(1/4,1/2)$ be fixed. For all
sufficiently large $N$ of either parity, there is a nonzero holomorphic
function $v_N(x,t)$ on a fixed complex neighborhood of $I$ and a fixed
color disk such that
\[
 v_N(r/N,t)=\frac{U_{Nt,r}}{U_{Nt,r+1}},\qquad
 v_N=\mu_s+N^{-1}\nu+O(N^{-2}).
\]
The identity is on the physical index grid, and the expansion holds
uniformly holomorphically. For all physical indices in $I$,
\begin{equation}\label{qi:high-index-leading}
 U_{Nt,r}=C(r/N,t)e^{NS(r/N,t)}(1+O(N^{-1})),
\end{equation}
with the $C,S$ of \eqref{qi:SC}. There is also a nonzero
holomorphic interpolation $H_N(x,t)$ on a smaller fixed domain with
relative grid error $O(e^{-cN})$ and the same leading expression.
All statements include both parities.
\end{lemma}
\begin{proof}
At $t=0$, throughout any compact subinterval of $(1/4,3/4)$,
\[
 \Re\mu_s>0,\qquad
 \left|\frac{\mu_f}{\mu_s}\right|=|\cot^2\pi x|<1.
\]
Choose $\eta>0$ with
$[\tau_0-\eta,1-\tau_0+4\eta]\subset(1/4,3/4)$.
The root labels extend to a fixed complex neighborhood and a
sufficiently small color disk, on which
\begin{equation}\label{qi:high-index-gap}
 \Re\mu_s\ge m>0,\qquad |\mu_f/\mu_s|\le k<1.
\end{equation}
The second root may vanish at $x=1/2$; the selected root does not.
Define $e=1/N$,
$\mathcal F_x(w)=p(x,t,e)+q_c(x,t,e)/w$, and
$L_N=\lceil(1/2-\tau_0+2\eta)N\rceil$. Set
\begin{equation}\label{qi:high-index-ratio}
 v_N(x,t)=\mathcal F_x\circ\mathcal F_{x+e}\circ\cdots
 \circ\mathcal F_{x+(L_N-1)e}
       \bigl(\mu_s(x+L_Ne,t)\bigr).
\end{equation}
We construct this function on a neighborhood of $[\tau_0,1/2]$,
and finally restrict it to $I$. All shifted arguments stay in the
larger neighborhood just chosen. One
backward step has defect $O(e)$ at the selected root and derivative
of modulus at most $k+O(e)$. Disks of radius $Be$ about the
successive roots are therefore preserved when $B$ is larger than
the uniform defect constant divided by $1-k$. They avoid zero.
Thus \eqref{qi:high-index-ratio} is holomorphic, nonzero and equal
to $\mu_s+O(e)$.

For a physical $x=r/N$, this composition passes through the row
$M-1$. There $q_c=0$, so all dependence on the artificial terminal value is erased and
$v_N((M-1)/N,t)=p((M-1)/N,t)$. The wall values
\eqref{qi:wall} and backward substitution prove the ratio identity.
They also prove nonvanishing of every $U_{Nt,r}$ used in the product;
there is no division by an unknown zero of the source.

Put $w=\mu_s+e\nu$. Expanding
$w(x)-\mathcal F_x(w(x+e))$ shows that its coefficient of $e$
vanishes precisely for the $\nu$ in \eqref{qi:SC}.
Its remaining defect is $O(e^2)$. The same contraction shows
$v_N=w+O(e^2)$, up to a terminal error bounded by $O(e)k^{L_N}$,
which is exponentially small. All bounds are holomorphic, so
Cauchy's estimate gives every fixed derivative bound on a smaller
domain.

The source is the finite product
\[
 U_{Nt,r}=\prod_{j=r}^{M-1}v_N(j/N,t).
\]
Let $f_N=\Log v_N$, with the selected logarithm, and let $b_N=M/N$.
Apply the Euler primitive \eqref{qi:EM} with $l_N=f_N$ and upper
endpoint $b_N$. This is a formula for the finite product just
displayed, independent of any transport argument. Cauchy's bound
$|\partial_x^j f_N|\le B_0j!a^{-j}$ and the periodic Bernoulli
remainder bound the error after $K_N=\lfloor c_0N\rfloor$ terms by
\[
 C N\left(\frac{K_N}{\pi aN}\right)^{2K_N}
       \le C'e^{-cN}
\]
for a sufficiently small fixed $c_0>0$. Exponentiating that explicit
primitive defines $H_N$ and proves its exponentially accurate grid
identity. The same bounds on a smaller complex tube show that the
sum of its Bernoulli terms beyond the integral and half-endpoints
is $O(N^{-1})$.

Since $f_N=\Log\mu_s+e\nu/\mu_s+O(e^2)$, its leading expression is
\[
 \sqrt{\mu_s(x,t)}
 \exp\left\{N\int_x^{1/2}\Log\mu_s(u,t)\dd u
          +\int_x^{1/2}\frac{\nu(u,t)}{\mu_s(u,t)}\dd u\right\}
 (1+O(N^{-1})).
\]
Here $\mu_s(1/2,t)=1$ and $b_N-1/2=O(N^{-1})$; the missing
integral of $\Log\mu_s$ is $O(N^{-2})$, so it changes the
exponent by only $O(N^{-1})$ at either parity. This proves
\eqref{qi:high-index-leading} with exactly the endpoint normalization
in \eqref{qi:SC}.
\end{proof}

At zero color, direct substitution in the rational coefficients gives
\[
 \Log\mu_s(x,0)=\pi\ii-2\pi\ii x,\qquad
 \frac{\nu(x,0)}{\mu_s(x,0)}
       =-\frac\pi2\cot(\pi x)-\frac{5\pi\ii}{2}.
\]
Integrating from the same upper endpoint yields
\begin{equation}\label{qi:high-index-complete}
 S(x,0)=\pi\ii(1/2-x)^2,\qquad
 C(x,0)=\sqrt{\sin\pi x}\,
       e^{-3\pi\ii(1/2-x)/2}.
\end{equation}
